% ============================================================================
% Six-Layer Strategic Allocation Framework Architecture Diagram
% Publication-quality TikZ diagram with professional styling
% ============================================================================

\begin{figure}[h]
\centering
\begin{tikzpicture}[
    % Professional color scheme
    layer1/.style={rectangle, draw=none, fill=blue!10, rounded corners=3pt, minimum width=14.5cm, minimum height=1.8cm, align=center, font=\small},
    layer2/.style={rectangle, draw=none, fill=green!10, rounded corners=3pt, minimum width=14.5cm, minimum height=1.8cm, align=center, font=\small},
    layer3/.style={rectangle, draw=none, fill=orange!10, rounded corners=3pt, minimum width=14.5cm, minimum height=1.8cm, align=center, font=\small},
    layer4/.style={rectangle, draw=none, fill=red!10, rounded corners=3pt, minimum width=14.5cm, minimum height=1.8cm, align=center, font=\small},
    layer5/.style={rectangle, draw=none, fill=purple!10, rounded corners=3pt, minimum width=14.5cm, minimum height=1.8cm, align=center, font=\small},
    layer6/.style={rectangle, draw=none, fill=teal!10, rounded corners=3pt, minimum width=14.5cm, minimum height=1.8cm, align=center, font=\small},

    % Component boxes
    component/.style={rectangle, draw=black!60, thick, rounded corners=2pt, minimum width=3.2cm, minimum height=1.2cm, align=center, font=\footnotesize\bfseries},

    % Arrows
    arrow/.style={->, >=stealth, very thick, color=black!70},
    dashedarrow/.style={->, >=stealth, thick, color=black!50, dashed},

    % Labels
    label/.style={font=\scriptsize\itshape, color=black!60}
]

% Layer 1: Protection Definition (AHP + Fuzzy Evaluation)
\node[layer1] (L1) at (0, 8) {};
\node[font=\bfseries\small] at (0, 8.6) {Layer 1: Protection Definition \& Priority Evaluation};
\node[component, fill=blue!20, minimum width=2.8cm] (ahp) at (-5.5, 8) {AHP\\Weighting};
\node[component, fill=blue!20, minimum width=2.8cm] (fuzzy) at (-1.8, 8) {Fuzzy\\Evaluation};
\node[component, fill=blue!20, minimum width=2.8cm] (species) at (1.9, 8) {Species\\Priority};
\node[component, fill=blue!20, minimum width=2.8cm] (gap) at (5.6, 8) {Protection\\Gap};
\node[label] at (0, 7.55) {Outputs: Protection scores $P_i$, species weights $\pi_s$, gaps $G_i$};

% Layer 2: Spatial Representation (Graph Theory)
\node[layer2] (L2) at (0, 5.8) {};
\node[font=\bfseries\small] at (0, 6.4) {Layer 2: Spatial Representation \& Network Model};
\node[component, fill=green!20, minimum width=4.2cm] (graph) at (-3.7, 5.8) {Graph Model\\$G=(V,E,W)$};
\node[component, fill=green!20, minimum width=4.2cm] (access) at (1.2, 5.8) {Accessibility \&\\Response Time};
\node[component, fill=green!20, minimum width=4.2cm] (zones) at (6.1, 5.8) {50 Monitoring\\Zones};
\node[label] at (0, 5.35) {Outputs: Network topology, accessibility matrix $A_{ik}$, travel times $T_{response}$};

% Layer 3: Resource Allocation (Integer Programming - MILP Core)
\node[layer3] (L3) at (0, 3.6) {};
\node[font=\bfseries\small] at (0, 4.2) {Layer 3: Resource Allocation \& Planning (MILP Core)};
\node[component, fill=orange!20, minimum width=4.2cm] (milp) at (-3.7, 3.6) {Integer Programming\\Optimization};
\node[component, fill=orange!20, minimum width=4.2cm] (constraint) at (1.2, 3.6) {Budget, Personnel,\\Coverage Constraints};
\node[component, fill=orange!20, minimum width=4.2cm] (multiobj) at (6.1, 3.6) {Multi-Objective\\Maximization};
\node[label] at (0, 3.15) {Outputs: Optimal deployment $\{x_{ijk}, y_{ij}, z_{ik}\}$, resource allocation patterns};

% Layer 4: Temporal Evaluation (Dynamic Programming + Queueing)
\node[layer4] (L4) at (0, 1.4) {};
\node[font=\bfseries\small] at (0, 2.0) {Layer 4: Temporal Evaluation \& Staffing Inference};
\node[component, fill=red!20, minimum width=4.2cm] (dp) at (-3.7, 1.4) {Dynamic Programming\\(Bellman Equation)};
\node[component, fill=red!20, minimum width=4.2cm] (queue) at (1.2, 1.4) {Queueing Model\\(M/M/c)};
\node[component, fill=red!20, minimum width=4.2cm] (staffing) at (6.1, 1.4) {Staffing\\Inference};
\node[label] at (0, 0.95) {Outputs: Temporal deployment policy, minimum staffing $c_{min}$, response times $W_q$};

% Layer 5: Uncertainty Handling (Monte Carlo Simulation)
\node[layer5] (L5) at (0, -0.8) {};
\node[font=\bfseries\small] at (0, -0.2) {Layer 5: Robustness \& Uncertainty Analysis};
\node[component, fill=purple!20, minimum width=4.2cm] (monte) at (-3.7, -0.8) {Monte Carlo\\Simulation};
\node[component, fill=purple!20, minimum width=4.2cm] (scenario) at (1.2, -0.8) {Scenario \&\\Sensitivity Analysis};
\node[component, fill=purple!20, minimum width=4.2cm] (ci) at (6.1, -0.8) {Confidence\\Intervals};
\node[label] at (0, -1.25) {Outputs: Robustness metrics, 95\% CI, probability distributions (10,000 trials)};

% Layer 6: Game-Theoretic Enhancement (Stackelberg Deterrence)
\node[layer6] (L6) at (0, -3.0) {};
\node[font=\bfseries\small] at (0, -2.4) {Layer 6: Adversarial Reasoning \& Deterrence (Optional)};
\node[component, fill=teal!20, minimum width=4.2cm] (stackelberg) at (-3.7, -3.0) {Stackelberg Game\\(Leader-Follower)};
\node[component, fill=teal!20, minimum width=4.2cm] (deterrence) at (1.2, -3.0) {Deterrence\\Quantification};
\node[component, fill=teal!20, minimum width=4.2cm] (nash) at (6.1, -3.0) {Nash\\Equilibrium};
\node[label] at (0, -3.45) {Outputs: Deterrence effects, adversarial best-response, strategy robustness};

% Vertical data flow arrows (main pipeline)
\draw[arrow] (0, 7.1) -- (0, 6.55);  % L1 -> L2
\draw[arrow] (0, 5.25) -- (0, 4.65); % L2 -> L3
\draw[arrow] (0, 3.05) -- (0, 2.45); % L3 -> L4
\draw[arrow] (0, 1.85) -- (0, 1.25); % L4 -> L5
\draw[arrow] (0, 0.35) -- (0, -0.25); % L5 -> L6

% Feedback arrows (dashed, showing iterative refinement)
\draw[dashedarrow] (-6.5, -2.6) -- (-6.5, 8.2);  % L6 -> L1 feedback
\draw[dashedarrow] (6.5, -2.6) -- (6.5, 4.2);   % L6 -> L3 feedback

% Feedback labels
\node[font=\tiny\itshape, rotate=90] at (-6.8, 2.5) {Parameter Calibration};
\node[font=\tiny\itshape, rotate=90] at (6.8, 0.5) {Strategy Refinement};

% Side annotation box
\node[rectangle, draw=black!40, dashed, rounded corners, align=left, font=\scriptsize, inner sep=4pt]
    at (-8.5, 2.5) {
        \textbf{Key Methods:}\\
        • AHP-fuzzy evaluation\\
        • Graph theory\\
        • Integer programming\\
        • Dynamic programming\\
        • Queueing theory\\
        • Monte Carlo\\
        • Game theory\\
        • Genetic algorithm\\
        (solver)
    };

% Final output box
\node[rectangle, draw=black!80, very thick, fill=yellow!10, rounded corners=5pt,
      minimum width=14.5cm, minimum height=1.2cm, align=center, font=\small\bfseries]
    at (0, -5.0) {
        FINAL OUTPUT: Strategic Resource Allocation Plan\\
        \font\small\mdseries 127 rangers, 89\% species protection, 94\% area coverage, 81\% threat detection
    };

% Arrow to final output
\draw[arrow, ultra thick] (0, -3.9) -- (0, -4.3);

\end{tikzpicture}
\caption{Six-layer strategic allocation framework for wildlife protection under uncertainty}
\label{fig:framework}
\end{figure}
